\section{Hamilton--Jacobi and Madelung Theory}
The Hamilton--Jacobi equation is usually presented as a
nonlinear first-order partial differential equation for a scalar field
$S(q,t)$. However, the fundamental object of classical mechanics is not a
one-point function but the two-point classical action $S(q_0,t_0;q,t)$, defined
as the extremal action between endpoints. The apparent reduction from a
two-point object to a one-point field hides essential geometric and dynamical
information. In particular, the Hamilton--Jacobi formulation is valid only
while the evolving Lagrangian manifold remains a graph over configuration
space. After caustic formation, the velocity field ceases to be globally
representable as a gradient, and the single-valued phase function breaks down.
We analyze the implications of this fact for Hamilton--Jacobi theory, Madelung
hydrodynamics, and stochastic formulations of quantum mechanics.

\subsection{Introduction}

The Hamilton--Jacobi (HJ) equation occupies a central place in classical
mechanics, geometrical optics, semiclassical quantum mechanics, and stochastic
formulations of quantum theory. It is usually written as
\begin{equation}
\partial_t S + H(q,\nabla S,t)=0,
\label{HJ}
\end{equation}
where $S(q,t)$ is interpreted as a scalar field on configuration spacetime.

This formulation suggests that the theory is an ordinary initial value problem
for a scalar field. However, this interpretation is misleading. The fundamental
classical object is not the one-point function $S(q,t)$ but the two-point
action $S(q_0,t_0;q,t)$, defined as the extremal value of the action functional
between fixed endpoints.

The purpose of this article is to clarify:
\begin{enumerate}
\item how the HJ equation emerges from the two-point action;
\item why the apparent one-point formulation hides essential boundary
    information;
\item why the velocity field ceases to be globally representable as a gradient
    after caustic formation;
\item why Madelung hydrodynamics and related stochastic formulations inherit
    these difficulties.
\end{enumerate}

\subsection{The Classical Action as a Two-Point Function}

Consider a classical system with Lagrangian
$L(q,\dot q,t)$.
The Hamilton principal function is defined by
\begin{equation}
S(q_0,t_0;q,t)=\operatorname*{ext}_{q(\tau)}
\int_{t_0}^{t} L(q,\dot q,\tau)\,d\tau,
\label{principal}
\end{equation}
where the extremization is performed over all paths satisfying
$q(t_0)=q_0$, $q(t)=q$.
Thus the classical action is naturally a two-point function on configuration spacetime.

The endpoint variation formulas yield
\begin{equation}
p_i(t)=\frac{\partial S}{\partial q^i},
\qquad
p_i(t_0)=-\frac{\partial S}{\partial q_0^i}.
\label{endpoint}
\end{equation}
Hence the two-point action generates the classical canonical transformation.

\subsection{Emergence of the Hamilton--Jacobi Equation}

Fix the initial endpoint $(q_0,t_0)$ and define
$\mathcal S(q,t):=S(q_0,t_0;q,t)$.
Now vary only the final endpoint. A small variation gives
$\delta S=p_i\,\delta q^i - H\,\delta t$.
Therefore,
$dS = p_i\,dq^i - H\,dt$.
Identifying $p_i=\frac{\partial S}{\partial q^i}$,
one obtains $\frac{\partial S}{\partial t}+
H\left(q,\frac{\partial S}{\partial q},t\right)=0$.
Thus the HJ equation appears only after one endpoint has been frozen.

\subsection{Hidden Boundary Data}

The reduction from the two-point action to the one-point field hides the
initial endpoint inside the boundary conditions.
Indeed, $S(q,t)=S(q_0,t_0;q,t)$
must satisfy singular initial data:
\begin{equation}
S(q,t_0)=
\begin{cases}
0, & q=q_0,\\
+\infty, & q\neq q_0.
\end{cases}
\end{equation}
Therefore the HJ equation alone does not define the dynamics. The choice of
admissible solutions depends on hidden global information inherited from the
original two-point action.

\subsection{Lagrangian Manifolds}

The geometric meaning of the HJ equation becomes clearer in phase space.
Given a function $S(q,t)$, define $p=\nabla S(q,t)$.
This determines a submanifold $\Lambda_t=\{(q,p):p=\nabla S(q,t)\}\subset T^*Q$.

Initially, one has a graph $\Lambda_0=\{(q,\nabla S_0(q))\}$.
Hamiltonian evolution transports this manifold: $\Lambda_t=\Phi_t(\Lambda_0)$,
where $\Phi_t$ is the Hamiltonian flow.The HJ formulation remains valid only 
while $\Lambda_t$ remains a graph over configuration space.

\subsection{Caustics and Breakdown of the Gradient Representation}

As the Hamiltonian flow evolves, the Lagrangian manifold generically folds in
phase space. The projection $\pi:\Lambda_t\to Q$ eventually ceases to be
one-to-one. At such points:
\begin{itemize}
\item several classical trajectories reach the same configuration point;
\item several momenta correspond to the same position;
\item the action becomes multivalued;
\item no single velocity field exists.
\end{itemize}
Instead of a single branch $p=\nabla S(q,t)$,
one obtains several branches: $p_1(q,t),\quad p_2(q,t),\quad \dots$
Each branch may locally be represented as a gradient,
$p_i=\nabla S_i$, but globally no single-valued $S(q,t)$ exists.

\subsubsection{Physical Interpretation}

Suppose particles initially move to the right. A sufficiently strong potential
may reverse some trajectories in finite time. Then at later times a single
position may simultaneously contain:
\begin{itemize}
\item trajectories moving left;
\item trajectories moving right.
\end{itemize}
Therefore the velocity field is no longer single-valued.
This phenomenon is generic and corresponds mathematically to the intersection
of characteristics in first-order PDE theory.

\subsection{Relation to Burgers Equation}

The HJ equation is closely related to Burgers-type equations.
For the free particle Hamiltonian $H=\frac{p^2}{2m}$,
one has $\partial_t S + \frac{(\nabla S)^2}{2m}=0$.
Defining $v=\frac{\nabla S}{m}$, gives $\partial_t v + (v\cdot\nabla)v=0$.
This equation develops shocks through characteristic crossing. Thus the
breakdown of the gradient representation is not exceptional but generic.

